\documentclass[11pt,a4paper]{article}

\usepackage[margin=24mm]{geometry}
\usepackage[T1]{fontenc}
\usepackage[utf8]{inputenc}
\usepackage{lmodern}
\usepackage{microtype}
\usepackage{xcolor}
\usepackage{booktabs}
\usepackage{tabularx}
\usepackage{array}
\usepackage{enumitem}
\usepackage{xurl}
\usepackage[hidelinks]{hyperref}
\emergencystretch=2em

\definecolor{pilotblue}{HTML}{2368FF}
\definecolor{pilotgreen}{HTML}{16834A}
\definecolor{pilotamber}{HTML}{B56A00}
\definecolor{pilotred}{HTML}{C52F46}
\newcommand{\Pilot}{\textbf{Pilot}}
\newcommand{\AION}{\textbf{AION Device Fabric}}
\newcommand{\decision}[1]{\par\medskip\noindent\colorbox{pilotblue!9}{\parbox{0.95\linewidth}{\textbf{Decided approach.} #1}}\par\medskip}

\title{\textbf{Pilot}\\Television-Led Adoption, Live Intelligence, Moments, and Agent Network Strategy}
\author{Tessaris / AION Product and Technical Record}
\date{29 August 2026\\Strategy decision following Fabric version 0.14.0}

\begin{document}
\maketitle

\begin{abstract}
This document records the decided product and distribution strategy for
Pilot, the consumer-facing intelligence operating on the AION Device Fabric.
Pilot will not initially be marketed as a general personal agent, household
operating system, or abstract artificial-intelligence platform.  Its entry
product is a highly demonstrable proposition: add conversational intelligence
to the television a person already owns, without requiring a new television or
dedicated AI appliance.  Entertainment, games, learning, search, live screen
intelligence, and moment sharing provide acquisition and immediate delight.
Once installed, the same governed mother brain becomes useful for reminders,
messages, calendars, shopping, booking, navigation, device control, safety,
and person-to-person agent handoffs.  The television is therefore the
adoption wedge; Pilot Actions creates habit; Pilot Relay creates network
growth; cross-device continuity creates retention.
\end{abstract}

\section{Status of this Record}

This is a product decision and implementation direction, not a claim that all
described capabilities are already complete.  The current Fabric proof has
already demonstrated discovery and governed control of an existing LG webOS
television, a mother node, a private phone controller, local voice activation,
television search and navigation, cloud-gaming composition, an interactive
learning centre, household persona boundaries, capability capsules, compact
device updates, and an auditable action chain.  Pilot Live, Pilot Moments,
cross-household Relay, car continuation, and emergency escalation are the next
strategic layers and require the additional work identified below.

\decision{The name \emph{Pilot} is the consumer identity and spoken activation
word.  AION Device Fabric remains the underlying technical substrate,
governance system, node architecture, memory, and execution boundary.}

\section{Core Product Thesis}

Consumers are unlikely to seek out and install a new ``personal autonomous AI
agent'' merely because it exists.  That description is abstract and asks for
trust before demonstrating value.  By contrast, the proposition

\begin{center}
\Large\textbf{``Add Pilot to the television you already own.''}
\end{center}

is concrete, visual, inexpensive to understand, and easy to demonstrate to
another person.  The user first experiences Pilot as an extraordinary upgrade
to a familiar device.  The broader personal agent emerges through use rather
than through an onboarding lecture.

\decision{AI television is the primary acquisition product.  It is not a
distraction from the general agent: it is the mechanism by which the general
agent enters the home, earns trust, pairs with the phone, learns permitted
preferences, and becomes available for wider action.}

The product should therefore be positioned in the following order:

\begin{enumerate}[leftmargin=*]
  \item \textbf{Install:} give an existing television Pilot capabilities.
  \item \textbf{Delight:} demonstrate entertainment control, games, learning,
        live understanding, search, and social sharing.
  \item \textbf{Habit:} let the user create reminders, lists, messages,
        calendar actions, research tasks, and shopping requests naturally.
  \item \textbf{Expansion:} continue the same identity, memory, and objective
        on the phone, computer, car, and other authorised nodes.
  \item \textbf{Network:} let useful moments and structured intents move to
        other people and their Pilots.
  \item \textbf{Retention:} preserve trusted context, preferences, completed
        work, device relationships, and verified service connections.
\end{enumerate}

\section{The Adoption Funnel}

\begin{table}[ht]
\centering
\small
\begin{tabularx}{\textwidth}{@{}>{\bfseries}p{19mm}XX@{}}
\toprule
Stage & What the user believes they installed & What Pilot is becoming\\
\midrule
Hook & AI for an existing television & A governed mother brain inside the home\\
Delight & Entertainment, games, learning, search, and live intelligence & Screen context, service knowledge, and useful memory\\
Habit & Reminders, lists, messaging, comparison, and scheduling & A personal action agent\\
Expansion & Phone, computer, car, calendar, and connected devices & One continuous intelligence across surfaces\\
Network & Sharing moments and requests with other people & A permissioned person-to-person agent network\\
Retention & ``Pilot already understands how I work'' & A persistent, trusted intelligence layer\\
\bottomrule
\end{tabularx}
\end{table}

The strategic relationship between the principal engines is:

\begin{center}
\textbf{Pilot TV acquires users} $\rightarrow$
\textbf{Pilot Live creates wonder} $\rightarrow$
\textbf{Pilot Moments spreads} $\rightarrow$\\[4pt]
\textbf{Pilot Actions creates daily use} $\rightarrow$
\textbf{Pilot Relay creates network effects} $\rightarrow$
\textbf{Pilot Everywhere creates permanence}.
\end{center}

\section{Pilot Live: Intelligence for Anything on Screen}

Pilot Live is the immediate room-scale intelligence layer.  It should
understand the current programme, application, conversation, event, and
available metadata sufficiently to answer and act on requests such as:

\begin{itemize}[leftmargin=*]
  \item ``Pilot, fact-check that.''
  \item ``Pilot, who is telling the truth here?''
  \item ``Pilot, explain what just happened without spoilers.''
  \item ``Pilot, translate what she just said.''
  \item ``Pilot, who is that actor?''
  \item ``Pilot, find the original source.''
  \item ``Pilot, explain that referee decision.''
  \item ``Pilot, teach the children the science behind this scene.''
  \item ``Pilot, is that advertised product actually any good?''
  \item ``Pilot, remember this; we should discuss it later.''
\end{itemize}

The differentiation is not the answer alone.  Pilot can move the answer into
memory, a message, a task, a shopping comparison, an educational activity, a
calendar item, a phone, a car, or a handoff to another person.

\subsection{Flagship fact-check proof}

The first headline demonstration should be:

\begin{center}
\Large\textbf{``We made an ordinary television fact-check a live broadcast
when you speak to it.''}
\end{center}

The intended flow is:

\begin{enumerate}[leftmargin=*]
  \item A factual statement is made in a live or recorded programme.
  \item The viewer says, ``Pilot, fact-check that.''
  \item Pilot identifies the relevant recent statement and extracts its
        falsifiable claims.
  \item The mother brain searches current, attributable evidence, preferring
        primary sources and clearly separating fact, inference, uncertainty,
        and disagreement.
  \item A compact television overlay displays a plain-language finding,
        confidence, and the most important comparison without unnecessarily
        replacing the programme.
  \item The user can ask a contextual follow-up such as ``show me the
        evidence'' without repeating the activation word during the bounded
        dialogue window.
  \item Full sources and private detail transfer to the paired phone.
\end{enumerate}

Because restricted LG televisions do not expose a supported general-purpose
screen-capture interface to third-party applications, the initial proof should
fuse bounded evidence rather than claim perfect visual access:

\begin{itemize}[leftmargin=*]
  \item a short, ephemeral, memory-only audio ring buffer;
  \item local speech transcription triggered for use only by an explicit
        request;
  \item foreground application, media, and programme-guide metadata;
  \item streaming-service title and playback-position evidence where exposed;
  \item phone-camera observation when visual evidence is necessary;
  \item claim extraction, current research, source grading, and contradiction
        detection; and
  \item a sanitized TV overlay plus a private evidence view on the phone.
\end{itemize}

Audio in the rolling buffer must be overwritten continuously and must not be
persisted by default.  An explicit user action is required before derived
transcript content is retained as part of a request or audit receipt.

\section{Pilot Moments: Entertainment as the Viral Loop}

Pilot Moments turns something seen on television into a social object.  The
canonical demonstration is:

\begin{quote}
``Pilot, share the last thirty seconds with Mike.''\\
``I found the moment.  Send it to Mike?''\\
``Yes.''
\end{quote}

If Mike already has Pilot, the moment can open on his phone or television and
retain its conversational context.  If Mike does not have Pilot, he receives a
normal secure link or provider-supported share.  The landing surface delivers
the value first and only then offers \emph{Watch on TV with Pilot}.  The content
therefore sells the product more effectively than an abstract invitation to an
AI platform.

\subsection{Service-aware sharing hierarchy}

Pilot must not bypass copyright controls, DRM, service terms, or access
controls.  ``Clip that'' is a user intent whose permitted execution depends on
the active source:

\begin{enumerate}[leftmargin=*]
  \item \textbf{Provider-native clip.} Invoke an official clipping function
        when the creator and service permit it.  Twitch, for example, exposes
        viewer clip creation and authorised sharing.\cite{twitchclips}
  \item \textbf{Provider-native moment.} Invoke an authorised moment-saving
        flow such as Netflix Moments when it is exposed to the signed-in user
        and available on the relevant device or service surface.\cite{netflixmoments}
  \item \textbf{Timestamp link.} Share the official title at the precise
        playback position where the service supports deep linking.  YouTube is
        moving toward Share at Timestamp as its principal viewer-sharing
        mechanism.\cite{youtubetimestamp}
  \item \textbf{Pilot Moment Card.} When video export is unavailable, create a
        copyright-safe card containing the title, episode, timestamp, official
        watch link, user reaction, short contextual description, and an
        ``open at this moment'' action where supported.
  \item \textbf{Media export.} Create a video file only for user-owned,
        public-domain, explicitly licensed, or provider-authorised content.
\end{enumerate}

Protected streams are intentionally encrypted and access-controlled; Pilot
must integrate with their approved interfaces rather than attempting to
defeat those protections.\cite{fairplay}

\subsection{Moment object}

A normalized moment record should contain no more information than necessary:

\begin{verbatim}
moment_id          signed opaque identifier
sender_persona     public/contact identity, not account credentials
service            netflix | youtube | twitch | broadcast | other
content_reference  provider title identifier or programme identifier
position           timestamp or bounded live-window reference
share_mode         native_clip | native_moment | timestamp | moment_card
reaction           sender-approved text or voice-derived caption
recipient          selected Pilot contact or external share destination
permissions        expiry, replay, forwarding, and audience policy
evidence_hash      hash of the normalized context used to prepare the share
\end{verbatim}

The card must not reveal a household member's viewing history, service token,
profile identity, private transcript, location, or unrelated context.

\section{Pilot Actions: The Personal Agent Revealed Through Use}

After installation for television intelligence, the user naturally begins to
delegate ordinary work:

\begin{itemize}[leftmargin=*]
  \item ``Pilot, add that to my to-do list.''
  \item ``Pilot, remind me tomorrow.''
  \item ``Pilot, message John that I will be twenty minutes late.''
  \item ``Pilot, reply to that email and say Thursday works.''
  \item ``Pilot, put this in my calendar.''
  \item ``Pilot, order more dog food.''
  \item ``Pilot, find the best three options and send them to my phone.''
  \item ``Pilot, arrange someone to repair the boiler.''
  \item ``Pilot, move the delivery to Friday.''
  \item ``Pilot, send this destination to the car.''
\end{itemize}

Shopping comparison is therefore a capability of Pilot Actions rather than a
separate strategic product.  The durable execution pattern remains:

\begin{center}
\textbf{Understand $\rightarrow$ research $\rightarrow$ prepare exact action
$\rightarrow$ private approval $\rightarrow$ execute once $\rightarrow$
verify $\rightarrow$ receipt.}
\end{center}

Shared television speech must never silently choose a person's card, account,
calendar, recipient, or private identity.  Consequential actions transfer to
the relevant person's private phone and connected service adapter.

\section{Pilot Relay: Person-to-Person Agent Handoffs}

Pilot Relay is the distribution engine.  Its central insight is:

\begin{center}
\large\textbf{The sender creates the recipient's first useful Pilot task before
the recipient has installed Pilot.}
\end{center}

For example, one person says:

\begin{quote}
``Pilot, ask Kevin to collect the dry cleaning on his way home.''
\end{quote}

The sender's Pilot creates a signed proposal, not a remote command.  Kevin's
Pilot may respond privately:

\begin{quote}
``Becca asked you to collect the dry cleaning on your way home.  Shall I remind
you when you leave work and add it to your route?''
\end{quote}

If Kevin accepts, his own Pilot---not Becca's---chooses an appropriate reminder
time, checks closing hours, and offers the stop to the navigation surface.  On
completion, it asks whether Kevin wants to send an acknowledgement.  Becca
does not gain access to Kevin's location, calendar, route, contacts, or
messages.

\subsection{Structured intent rather than unstructured messaging}

Relay should transport a typed, bounded proposal:

\begin{verbatim}
type:       task
from:       sender contact identity
to:         recipient contact identity
action:     collect
object:     dry cleaning
window:     journey home
location:   recipient-resolved usual dry cleaner
status:     proposed
return:     optional completion acknowledgement
\end{verbatim}

Additional handoff types include message, event, invitation, recommendation,
shopping item, location, document, payment request, decision, emergency alert,
introduction, moment, and shared plan.  Typed objects let the recipient's Pilot
reason about time, place, interruption, dependencies, and permitted execution
without conflating every request with ordinary chat.

\subsection{Installed and uninstalled recipients}

If the recipient has Pilot, the proposal is delivered through the recipient's
trusted contact and capability policy.  If the recipient does not have Pilot,
the system sends a conventional, sender-approved message containing a secure
temporary page.  The recipient can accept, add the request to an existing
phone tool, decline without creating an account, or activate Pilot.  Contact
discovery and SMS invitation are established low-friction patterns in existing
communications products.\cite{whatsappinvite}

Pilot must support households but must not be limited to a rigid household
account.  The useful graph includes partners, children, parents, friends,
colleagues, neighbours, carers, tradespeople, school contacts, and temporary
groups.  Conventional family-account systems are often deliberately small;
for example, Google documents a family manager inviting up to five other
people.\cite{googlefamily}  Pilot instead requires a permissioned contact graph
and temporary or persistent \emph{Circles} for household, trip, school, care,
work, and friend contexts.

\subsection{Non-negotiable trust rule}

\decision{One person's Pilot must never command another person's Pilot.  The
sender's Pilot proposes; the recipient's Pilot evaluates; the recipient
controls acceptance, scheduling, disclosure, execution, and acknowledgement.}

Recipient policies must support at least:

\begin{itemize}[leftmargin=*]
  \item always allow ordinary reminders from selected contacts;
  \item require confirmation before calendar modification;
  \item never share location or presence by default;
  \item silence non-urgent proposals during protected periods;
  \item allow explicitly configured emergency interruption;
  \item block, mute, expire, or rate-limit a contact;
  \item forbid purchases or payments requested by another person;
  \item prevent private proposals from appearing on communal screens; and
  \item apply transparent, age-appropriate guardian policy for children and
        teenagers without creating invisible surveillance.
\end{itemize}

Every handoff requires signed sender identity, intended recipient, expiry,
nonce or replay protection, scope, audit receipt, and a human-readable preview.

\section{Pilot Companion, Guardian, Everywhere, and Work}

\subsection{Pilot Companion}

Pilot Companion provides contextual conversation rather than a detached
chatbot.  It may discuss the active programme, remember permitted preferences,
explain confusing material, conduct quizzes, continue a conversation on
another device, and help a user contact family.  It must not claim to be human,
manipulate emotional dependence, manufacture urgency, or interrupt without a
user-controlled policy.  Warmth and continuity are product qualities; honesty
and user agency are system requirements.

\subsection{Pilot Guardian}

The long-term safety proposition includes requests such as ``Pilot, help me'',
``Pilot, I have fallen'', and ``Pilot, call an ambulance''.  A safe proof may
acknowledge the call, display a large emergency interface, contact configured
trusted people, share explicitly authorised location, and open a two-way
response surface.  Automatic emergency-service escalation cannot be marketed
or relied upon until the system has jurisdiction-specific authorised calling,
reliable identity and location, connectivity fallback, false-positive
handling, health and emergency review, continuous availability monitoring,
and a tested failure protocol.  Guardian is safety infrastructure, not an
ordinary voice shortcut.

\subsection{Pilot Everywhere}

The intelligence must continue across surfaces rather than being copied into
unrelated assistants.  A restaurant shortlist created on the television can
move privately to the phone; the selected destination can later appear in the
car; a voice note from the car can become a work task; a household request can
become an authorised navigation stop.  The same persona, objective, memory,
permissions, and audit chain continue while each surface receives only the
capabilities and data appropriate to it.

\subsection{Pilot Work}

Work mode remains a high-value expansion: live briefings, research walls,
meeting preparation, document review, operational dashboards, scenario
analysis, action extraction, and governed execution.  It is deliberately not
the initial consumer acquisition message.  The television-led consumer proof
establishes the cross-device substrate on which a later room-scale work and
Boardroom product can operate.

\section{Fabric Architecture Required by the Strategy}

\begin{table}[ht]
\centering
\small
\begin{tabularx}{\textwidth}{@{}p{38mm}X@{}}
\toprule
Fabric component & Strategic responsibility\\
\midrule
Mother brain & Reasoning, planning, service orchestration, private memory,
persona resolution, policy, and recovery.\\
Universal node identity & Stable public identity and stored handshakes that
allow one device node to serve multiple independently authorised mothers.\\
Television gateway/node & Governed control, current-service evidence, Canvas
presentation, and verified restricted-device actions.\\
Private phone companion & Identity selection, private detail collection,
contact choice, camera observation, approval, and consequential-action receipt.\\
Pilot Live context layer & Ephemeral recent-audio context, metadata fusion,
claim extraction, current research, and evidence-labelled response.\\
Moment adapter layer & Provider-native clips, moments, timestamp links, and
copyright-safe fallback cards.\\
Relay envelope & Signed typed proposal with recipient, scope, expiry, replay
protection, privacy policy, status, and optional acknowledgement.\\
Service adapters & Calendar, email, messaging, tasks, maps, shopping, booking,
media, and later authorised emergency providers.\\
Audit and glyph ledgers & Compact, verifiable records of permission, action,
state transition, handoff, outcome, and reconstruction without raw secret data.\\
\bottomrule
\end{tabularx}
\end{table}

\section{Decided Implementation Sequence}

\subsection*{Phase 1: Pilot Live and Pilot Moments}

\begin{enumerate}[leftmargin=*]
  \item Add the ephemeral recent-audio context buffer and explicit activation
        boundary.
  \item Fuse current application, content, programme, playback-position, and
        phone-observer evidence into a normalized live-context object.
  \item Implement ``fact-check that'', ``explain that'', ``translate that'',
        and contextual follow-up on the television and phone.
  \item Implement the service-aware moment hierarchy: native share, timestamp,
        Moment Card, or explicit refusal.
  \item Add private contact selection and a single-recipient share receipt.
\end{enumerate}

\subsection*{Phase 2: Pilot Contacts and Relay}

Implement contact identity, signed invitations, installed-recipient delivery,
temporary value-first web delivery for uninstalled recipients, typed handoff
objects, recipient policies, expiry, replay protection, acknowledgement, and
block/mute controls.

\subsection*{Phase 3: Pilot Actions}

Connect real user-authorised task, messaging, calendar, email, maps, shopping,
and booking services to the existing prepare--approve--execute--verify flow.
Moment and Relay objects then become direct inputs to useful action.

\subsection*{Phase 4: Pilot Everywhere}

Provide phone background continuity, car-safe voice interaction, authorised
navigation handoff, location-triggered reminders, and multi-surface objective
resumption without exposing private state to unrelated nodes.

\subsection*{Phase 5: Pilot Guardian and Pilot Work}

Develop Guardian only with the necessary reliability, emergency, privacy, and
regulatory partners.  Expand Work and Boardroom mode using the proven identity,
handoff, service, and room-scale presentation layers.

\section{Success Measures}

The strategy should be measured as a behavioural system, not merely as a list
of available commands:

\begin{itemize}[leftmargin=*]
  \item time from download to first successful television action;
  \item percentage of activated households pairing a private phone;
  \item first-day use of a Live command or Moment share;
  \item successful provider-native or timestamp share rate;
  \item percentage of recipients receiving value before installation;
  \item recipient-to-Pilot activation conversion;
  \item accepted Relay proposals per active contact pair;
  \item weekly use of at least one Pilot Action outside entertainment;
  \item cross-surface continuation rate among television, phone, and car;
  \item false activation, failed execution, reversal, and complaint rates;
  \item percentage of consequential actions with explicit private approval and
        verified outcome; and
  \item user-controlled deletion, mute, block, and privacy-policy effectiveness.
\end{itemize}

Growth must not be optimized by dark patterns, involuntary contact uploads,
silent invitations, public exposure of private viewing, or pressure to accept
another person's request.  A network of personal agents is valuable only if
each person remains sovereign over their own agent.

\section{Final Product Position}

Pilot should not initially ask the market to understand an autonomous agent,
a device fabric, or a household operating system.  It should make one promise:

\begin{center}
\Large\textbf{Add Pilot to the television you already own.}
\end{center}

The demonstration then proves that Pilot can understand, fact-check, find,
remember, and share what the user is watching.  Entertainment creates the
installation; live intelligence creates the story people repeat; Moment
sharing creates the invitation; Actions creates daily utility; Relay turns
individual utility into a network; and continuity across phone, home, and car
makes the intelligence indispensable.

\begin{center}
\textbf{Screen $\rightarrow$ understanding $\rightarrow$ action
$\rightarrow$ another person.}
\end{center}

This is the decided approach: the television is the visible beginning, not the
product boundary.  Pilot enters through entertainment and becomes a governed,
persistent personal intelligence through demonstrated usefulness.

\begin{thebibliography}{9}

\bibitem{netflixmoments}
Netflix, ``Introducing Moments: Save, Relive and Share Your Favorite Netflix
Scenes,'' 28 October 2024.\\
\url{https://about.netflix.com/en/news/introducing-moments-save-relive-and-share-your-favorite-netflix-scenes}

\bibitem{twitchclips}
Twitch Help, ``How to Create, Edit, and Share Clips.''\\
\url{https://help.twitch.tv/s/article/how-to-use-clips}

\bibitem{youtubetimestamp}
TeamYouTube, ``Updating how you share video moments,'' 16 April 2026.\\
\url{https://support.google.com/youtube/thread/425735532}

\bibitem{fairplay}
Apple Developer, ``FairPlay Streaming.''\\
\url{https://developer.apple.com/streaming/fps/}

\bibitem{whatsappinvite}
WhatsApp Help Center, ``How to find and invite contacts.''\\
\url{https://faq.whatsapp.com/1183494482518500/}

\bibitem{googlefamily}
Google Account Help, ``Invite more members.''\\
\url{https://support.google.com/accounts/answer/15077834}

\end{thebibliography}

\end{document}
